% VI/49 reader bridge -- additive exposition, original exercises unchanged.
\paragraph{What an affine vanishing theorem actually says.}
If $X=\operatorname{Spec}A$ and $\mathcal F$ is a
\emph{quasi-coherent} $\mathcal O_X$-module, then
\[
H^i(X,\mathcal F)=0\qquad(i>0).
\]
The quasi-coherence hypothesis matters: this is not a
statement about every sheaf of abelian groups on every
affine scheme.  Its importance is that an affine chart
can be used as a cohomologically trivial computational
piece.  If $X$ is covered by affine opens whose finite
intersections are affine (for instance a separated
scheme with an affine cover), the \v{C}ech complex for a
quasi-coherent sheaf computes its sheaf cohomology.

\paragraph{Worked example 1: the affine line.}
Let $X=\operatorname{Spec}k[t]$ and
$\mathcal F=\widetilde{k[t]/(t^2)}$.  The global sections are
\[
H^0(X,\mathcal F)=k[t]/(t^2),
\qquad H^i(X,\mathcal F)=0\quad(i>0).
\]
The nilpotent class of $t$ remains visible in $H^0$:
vanishing of higher cohomology does not discard the
nonreduced structure.  The same result holds for the
sheaf associated with an arbitrary $k[t]$-module, even
one that is not finitely generated, because it is still
quasi-coherent.

\paragraph{Worked example 2: the exactness test over an affine base.}
For a ring $A$, every short exact sequence of
$A$-modules
$0\to M'\to M\to M''\to0$ gives
\[
0\to\widetilde M'\to\widetilde M\to
 \widetilde M''\to0
\]
on $\operatorname{Spec}A$.  Taking global sections
recovers the original module sequence because
$\Gamma(\operatorname{Spec}A,\widetilde M)=M$.
Consequently the connecting map
$H^0(\widetilde M'')\to H^1(\widetilde M')$ is zero.
This is a precise reason that quasi-coherent global
sections are exact on an affine scheme.

\paragraph{Worked example 3: one projective line, two affine charts.}
Let $X=\mathbb P^1_k$ with
$U_0=\operatorname{Spec}k[t]$,
$U_1=\operatorname{Spec}k[t^{-1}]$ and
$U_{01}=\operatorname{Spec}k[t,t^{-1}]$.
For $\mathcal O(n)$, all three opens are affine and
the sheaf remains quasi-coherent on restriction.  The
two-open \v{C}ech complex is therefore a valid
cohomology computation:
\[
0\longrightarrow
 k[t]\oplus k[t^{-1}]
 \xrightarrow{(a,b)\mapsto t^n b-a}
 k[t,t^{-1}]
 \longrightarrow0
\]
where the coefficients are expressed in the frame
$X_0^n$ on $U_0$.  The overlap quotient is
\[
H^1(X,\mathcal O(n))\simeq
\frac{k[t,t^{-1}]}{k[t]+t^n k[t^{-1}]}.
\]
Thus
\[
 h^0(\mathcal O(n))=\max\{n+1,0\},\qquad
 h^1(\mathcal O(n))=\max\{-n-1,0\}.
\]
For $n=-4$ a basis of $H^1$ is represented by
$t^{-1},t^{-2},t^{-3}$; hence $h^1=3$.
For $n=2$, the three global sections are
$X_0^2,X_0X_1,X_1^2$, with $h^1=0$.
In all integer degrees the Euler characteristic is
\[
 \chi(\mathbb P^1,\mathcal O(n))
 =h^0-h^1=n+1.
\]
This is the elementary genus-zero Riemann--Roch
calculation, derived here directly from Laurent powers.

\paragraph{What the calculation does \emph{not} imply.}
The scheme $\mathbb P^1$ is covered by two affine opens,
but it is not itself affine.  Its higher cohomology need
not vanish for all quasi-coherent sheaves:
$H^1(\mathbb P^1,\mathcal O(-2))\simeq k$ is a concrete
counterexample.  Conversely, the group
$H^1(\mathbb P^1,\mathcal O)$ vanishes; vanishing for
\emph{one} sheaf does not establish that a scheme is affine.

\paragraph{A cohomological-dimension bound from an affine cover.}
Suppose a separated scheme $X$ admits a cover by $r$
affine open subsets and $\mathcal F$ is quasi-coherent.
Every finite intersection of these opens is affine,
so the increasing-index \v{C}ech complex has no terms
in degree $r$ or higher.  Hence
\[
H^q(X,\mathcal F)=0\qquad(q\geq r).
\]
For $\mathbb P^1$ we may take $r=2$, which explains
why no $H^2$ term appears in the preceding calculation.
This simple bound is generally not optimal but follows
directly from the cover.

\paragraph{Flabby vanishing and the role of hypotheses.}
For a flabby sheaf $\mathcal G$ of abelian groups,
every restriction map $\mathcal G(U)\to\mathcal G(V)$
for $V\subseteq U$ is surjective; a flabby sheaf is
acyclic for sheaf cohomology:
\[
H^i(X,\mathcal G)=0\qquad(i>0).
\]
This theorem is independent of quasi-coherence, and
should not be confused with the affine theorem.  A
general sheaf may have higher cohomology even on an
affine scheme; quasi-coherent and flabby are two
different sufficient conditions in their respective
settings.

\paragraph{Positive twists: a statement of Serre vanishing.}
If $X$ is projective over a Noetherian ring, $\mathcal L$
is a relatively ample invertible sheaf, and
$\mathcal F$ is coherent, then for sufficiently large
integers $m$,
\[
 H^i(X,\mathcal F\otimes\mathcal L^{\otimes m})=0
 \qquad (i>0).
\]
In its standard form over an affine Noetherian base,
this is Serre vanishing.  The bound on $m$ depends on
the sheaf; it is not generally uniform over all
coherent sheaves.  For $\mathbb P^1_k$,
$\mathcal L=\mathcal O(1)$ and
$\mathcal F=\mathcal O(-4)$, we have
$\mathcal F\otimes\mathcal L^{\otimes m}
=\mathcal O(m-4)$, whose $H^1$ vanishes for $m\geq3$.
At $m=2$ the group is $H^1(\mathcal O(-2))\simeq k$,
so the explicit bound is sharp in this example.

\paragraph{Final connection for Part VI.}
The central computational chain is now complete:
module localization produces quasi-coherent sheaves;
affine charts suppress their higher cohomology;
overlap complexes detect global gluing obstructions;
long exact sequences transport those obstructions
between related sheaves; and positivity can make the
obstructions vanish.  These principles are the bridge
from the scheme and projective geometry chapters to
more advanced cohomological methods.
